% !TeX root = ./../main.tex

\chapter*{Streszczenie}
\thispagestyle{empty}
\agentinfo{
Streszczenie w języku polskim, około jednej strony.
Pisać na końcu, po ustabilizowaniu treści rozdziałów 4--8 i podsumowania.
Układ: problem (szum w erze NISQ a kwantowe przetwarzanie obrazów), cel i tezy, metody (PDU, LPIQE, pix2pix, modelowanie szumu), najważniejsze wyniki liczbowe, wnioski.
}

\noindent\textbf{Słowa kluczowe:} informatyka kwantowa, redukcja błędów kwantowych, korekcja błędów kwantowych, próbkowanie kwantowe, kwantowe przetwarzanie obrazów, klasyfikacja obrazów, era NISQ

\cleardoublepage

\begin{otherlanguage}{british}
\chapter*{Abstract}
\thispagestyle{empty}
\agentinfo{
Streszczenie w języku angielskim, wierny odpowiednik wersji polskiej.
Tytuł angielski pracy: makro TitleAlt w main.tex.
}

\noindent\textbf{Keywords:} quantum computing, quantum error reduction, quantum error correction, quantum sampling, quantum image processing, image classification, NISQ era
\end{otherlanguage}

\cleardoublepage
